\cotefile{156-7}
\scanpage{80}
\mspage{79}
% slanted
\margin{En dehors des cas \unsure{quasi-ensembliste|quasi-ensemblistes} et \unsure{ensembliste|ensemblistes} ces questions n'ont d'intérêt que pour des $A_i$ formés d'éléments mutuellement compatibles - - - -}

Qu'en est-il des formules pour les Sup et les Inf ? Moyennant passage aux complémentaires, on est réduit à celle sur les Inf, qui est l'inclusion
\[
\deleted{\unsure{\subset}}\quad |\bigcap A_i|\mathrel{\overset{\text{?}}{\overset{=}{\subset}}}\bigcap_i|A_i| \tag*{(X)}
\]
(les $A_i$ des parties-supports). \deleted{\unsure{Prenons} le complément des premiers, soit $\deleted{\unsure{U}}U$ il \deleted{\unsure{faut}} \deleted{\unsure{donc}} \unsure{montrer}}
\[
\deleted{U\cap\bigcap_i|A_i|=\varnothing\quad\text{??}}
\]
\deleted{Or par Hyp,}
\[
\deleted{U=|\bigvee_i\unsure{A'_i|A_i}|}
\]

Je ne sais pas aller plus loin, sauf si je fais l'hypothèse draconienne suivante, beaucoup plus forte que la précédente :
\[
\text{\uline{Hyp'}}\quad \forall a\in\unsure{\mathbb{R}|\mathbb{P}|P},\quad\exists X\in M\text{ avec }|X|=\{a\}.
\]
Ceci implique déjà la séparation $(\unsure{\mathrm{vi}'|\mathrm{vi}|\mathrm{v}'}_{\unsure{M|\mu}}\,1\mathrm{b}\unsure{,|)})$, la compatibilité de $\overline{\Sigma}_M\mathrel{\unsure{-|\longrightarrow}}\unsure{\mathcal{P}|\mathscr{P}}(P)$ aux applications $\mathrm{C}$ (plus gén. la formule
\[
|\operatorname{cosupp}A|=P\smallsetminus|A|\qquad\forall A\subset M,
\]
) je dis que ça implique aussi \uline{l'égalité dans} (X) ci dessus. Soit en effet un pt $a$ du second membre, \deleted{$\unsure{\forall_i}$}, soit $X\in M$ \inserted{\unsure{lui-même|équivalent}} avec $|X|=\{a\}$, on a \deleted{\unsure{donc} $\forall$} $|X|\subset|A_i|$ \deleted{$\unsure{\Rightarrow}$} $\operatorname{supp}X\subset\underset{\inserted{\unsure{=A_i|A_i}}}{\operatorname{supp}A_i}$
